\section{问题三模型的建立与求解}

\subsection{求解思路}
问题三在问题二的基础上增加每天 0:00、6:00、12:00、18:00 发布的未来 24 小时整点光伏预报（附件 3），
并允许在 6:00、12:00、18:00 依据新预报调整当天尚未执行的购电计划：少购部分按交易时刻电价的 50\% 结算，
多购部分按 150\% 结算，供电缺口仍按 5 倍电价紧急购电。因此关键不在于“用最新预报直接替换原计划”，
而在于回答\textbf{后续时刻的预报是否值得用于制定调整购电策略}。
本文先对附件 3 的预报质量作诊断，再建立“预报标定—当日残差修正—双向滚动调整—末端备用”的两阶段调度模型，
通过策略变体对比拆解“预报质量”与“调整机制”各自的贡献，定量回答上述问题。
求解流程如图~\ref{fig:q3-flow} 所示。

\begin{figure}[!htbp]
\centering
\footnotesize
\begin{tikzpicture}[
  >=Stealth,
  node distance=6mm,
  plan/.style={rectangle, draw=black!70, fill=orange!12, align=left, text width=140mm,
               minimum height=10mm, inner sep=3pt},
  stage/.style={rectangle, draw=black!70, fill=blue!6, align=left, text width=140mm,
               minimum height=10mm, inner sep=3pt},
  io/.style={rectangle, rounded corners=2mm, draw=black!70, fill=gray!12, align=center,
               text width=140mm, minimum height=8mm, inner sep=3pt},
  line/.style={draw, -Latex, thick}
]
  \node[plan] (p1) {\textbf{0:00　制定日前计划}：MOS 标定后的光伏估计与上周同星期实际值（w1）融合，
                    按预测余量 $\gamma=1.02$、$\delta=0.98$ 构造保守净负荷；\\
                    求解 $\min\sum_{t} c_t p^{plan}_t\Delta t$，
                    满足功率平衡、储能约束、日循环约束 $S^{plan}_{144}=S^{real}_{0}$};
  \node[stage] (a1) [below=of p1] {\textbf{6:00 / 12:00 / 18:00　双向滚动调整}：接收新预报
                    $\rightarrow$ MOS 标定 $+$ 当日已实现残差修正（$\lambda_{pv}=0.1$、$\lambda_L=0.3$）；\\
                    求解调整 LP：$\min\sum_{t\ge a}\left[c_t p^{adj}_t+0.5c_t(u^{+}_t+u^{-}_t)+5c_t e_t\right]\Delta t$，
                    日末最低备用 $S^{plan}_{144}\ge 4000$\,kWh};
  \node[stage] (r1) [below=of a1] {\textbf{全天 144 个 $10\,$min　日内实际运行}：实际负载与光伏
                    $\rightarrow$ 储能实时纠偏 $\rightarrow$ 缺口按 $5c_t$ 紧急购电、富余弃光；\\
                    按 $C(t)=c_t p^{adj}_t+0.5c_t\left|p^{adj}_t-p^{plan}_t\right|+5c_t e_t$ 结算};
  \node[io] (o1) [below=of r1] {全年购电总费用与四个指定日期的计划、调整结果
                    $\rightarrow$ 与问题二策略及完美预测下界的对比};

  \draw[line] (p1.south) -- node[right, font=\scriptsize] {计划购电量 $p^{plan}_t$} (a1.north);
  \draw[line] (a1.south) -- node[right, font=\scriptsize] {调整后购电量 $p^{adj}_t$} (r1.north);
  \draw[line] (r1.south) -- (o1.north);
\end{tikzpicture}
\caption{问题三“计划—调整—结算”两阶段滚动调度流程}
\label{fig:q3-flow}
\end{figure}

\subsection{模型建立}

\subsubsection{附件 3 预报质量诊断}
将附件 3 全年 365 天、每天 4 次发布的整点预报与附件 2 实际光伏出力按同一口径逐时对照
（“预报第 $k$ 小时”对应发布时刻后第 $k$ 个整点，实际值取该小时内 6 个 $10\,$min 区间均值），
统计各目标时刻的平均绝对误差（MAE）\cite{fpp3} 与系统性偏差，如图~\ref{fig:q3-forecast} 所示。

\begin{figure}[!htbp]
\centering
\includegraphics[width=0.485\textwidth]{Q3_mae.pdf}\hfill
\includegraphics[width=0.485\textwidth]{Q3_bias.pdf}
\caption{附件 3 预报质量诊断：（左）不同发布时刻的预报 MAE 随目标时刻的变化；
（右）预报偏差（预报值$-$实际值）随目标时刻的变化}
\label{fig:q3-forecast}
\end{figure}

由图~\ref{fig:q3-forecast}（左）可得两点结论：
\textbf{其一，同一目标时刻的预报精度几乎不随发布时间改善}——目标时刻 16:00 时，
0:00、6:00、12:00 发布的预报 MAE 分别为 1000.2、1043.6、1066.5\,kW，反而略有上升；
14:00 之后上周同星期实际值（w1）的 MAE 显著更低（16:00 时仅 214.7\,kW）。
\textbf{其二，“更新”的预报并不必然优于早前预报}——对目标时刻 12:00，
12:00 发布序列中对应的恰为 24 小时前作出的预报，其 MAE 高达 1374.1\,kW，
明显劣于当日 0:00 预报（640.9\,kW）与 6:00 预报（429.4\,kW）。

由图~\ref{fig:q3-forecast}（右）可知，预报存在明显的\textbf{系统性偏差}：
上午时段预报偏低（8:00 约 $-1050$\,kW，10:00 约 $-580\sim-700$\,kW），
午后时段预报偏高（14:00 约 $+567\sim+790$\,kW，15:00 约 $+836\sim+988$\,kW，16:00 约 $+988\sim+1067$\,kW），
且不同发布时刻的偏差方向高度一致，属预报本身的系统性偏差；
预报峰值时刻（13:00）与实际峰值时刻（12:00）亦存在错位。
此外，12:00 报的修订量与当日上午“实际$-$预报”偏差的相关系数仅 0.02$\sim$0.04，
说明\textbf{后续预报几乎不含当日天气的新信息}。

尽管信息增量有限，预报值与实际出力仍具有较高的相关性（0:00 报逐时相关系数 0.78 以上），
其系统性偏差可用历史数据标定修正；且当日上午实际出力与午后实际出力的相关系数达 0.914，
w1 残差的日内自相关系数光伏为 0.464、负载为 0.592，
说明\textbf{后续预报的正确用法是“经标定的输入 $+$ 当日已实现残差修正”，而非当作真值直接使用}。

\subsubsection{光伏出力估计模型}
采用模型输出统计（MOS）方法\cite{glahn1972} 对光伏出力标定，并以当日已实现残差修正后续时段。

\textbf{（1）0:00 计划用估计值。} 逐小时建立线性回归\cite{fpp3}，融合当日 0:00 官方预报与 w1：
\begin{equation}
\widehat{PV}(d,h)=a_h+b_h\,\widehat{PV}^{off}_{0:00}(d,h)+c_h\,\widehat{PV}^{w1}(d,h)
\label{eq:q3-mos}
\end{equation}
回归系数由最近 30 天的（官方预报，w1，实际值）三元组按最小二乘拟合并逐日滚动更新（保证因果性），
历史不足 5 天时退化为 0.5/0.5 等权融合。以 14:00 为例，
$\widehat{PV}\approx232+0.34\,\widehat{PV}^{off}_{0:00}+0.58\,\widehat{PV}^{w1}$，
MAE 为 373.0\,kW，优于单独使用 w1（456.8\,kW）与原始预报（793.4\,kW）。
逐时估计再按上周同日的日内形状系数分解到 144 个时段。

\textbf{（2）调整时刻的更新估计值。} 在调整时刻 $\tau\in\{6,12,18\}$，
将 $\tau$ 时刻发布的官方预报代入式~\eqref{eq:q3-mos} 得基准估计，并叠加当日已实现残差：
\begin{equation}
\widehat{PV}_\tau(h)=\widehat{PV}(h)+\lambda_{pv}\overline{\Big(PV^{real}-\widehat{PV}\Big)}_{h'<\tau},\qquad
\widehat{L}_\tau(h)=\widehat{L}(h)+\lambda_L\overline{\Big(L^{real}-\widehat{L}\Big)}_{h'<\tau}
\label{eq:q3-resid}
\end{equation}
其中 $\overline{(\cdot)}_{h'<\tau}$ 对 $\tau$ 之前已实现的整点时段取平均。
由于残差日内自相关为正但不强，$\lambda$ 宜取中小值；以全年总费用为目标网格寻优，
最终取 $\lambda_{pv}=0.1$、$\lambda_L=0.3$（见表~\ref{tab:q3-lambda}）。

\subsubsection{计划与调整两阶段优化模型}
\textbf{（1）费用结算口径。} 按题意，第 $t$ 个时段的购电费用为
\begin{equation}
C(t)=c_t\,p^{adj}_t+0.5\,c_t\left|p^{adj}_t-p^{plan}_t\right|+5\,c_t\,e_t
\label{eq:q3-cost}
\end{equation}
三项依次为按调整购电量实付的电费、相对计划的违约与溢价费用、紧急购电费用；
它精确等价于“计划内按 $c_t$、超出计划按 $1.5c_t$、少购部分按 $0.5c_t$”结算。
引入非负变量 $u^{+}_t,u^{-}_t$ 表示上调与下调量以便线性化：
\begin{equation}
p^{adj}_t-p^{plan}_t=u^{+}_t-u^{-}_t,\qquad u^{+}_t\ge0,\quad u^{-}_t\ge0 .
\label{eq:q3-split}
\end{equation}

\textbf{（2）0:00 计划模型。} 结构与问题二一致：余量取 $\gamma=1.02$、$\delta=0.98$
（调整机制提供了追索权，余量较问题二的 1.05 下降，寻优见第~\ref{sec:q3-param} 节），
以当日计划购电费用最小为目标，满足功率平衡、储能约束与日循环约束 $S_{144}=S_0$。

\textbf{（3）调整模型。} 在调整时刻 $\tau$ 对尚未执行的时段 $t\ge a$ 重新优化：
\begin{equation}
\min\ \sum_{t=a}^{144}\Big[c_t\,p^{adj}_t+0.5\,c_t\,u^{+}_t+0.5\,c_t\,u^{-}_t+5\,c_t\,e_t\Big]\Delta t
\label{eq:q3-adjlp}
\end{equation}
约束包括更新估计下的功率平衡、以该时刻实际储电量为初值的储能状态转移与容量、功率约束、
式~\eqref{eq:q3-split} 的调整量分解，以及紧急购电松弛变量 $e_t\ge0$。
$e_t$ 作为“安全阀”保证任意更新情形下模型可行，并使模型在“$1.5c_t$ 上调”与“$5c_t$ 紧急购电”之间自动比价。
此外，为避免模型在傍晚短视放空储能，设置\textbf{日末最低备用约束}
\begin{equation}
S^{plan}_{144}\ge S^{end}_{\min}=4000\ \text{kWh}
\label{eq:q3-end}
\end{equation}
该约束仅作用于调整计划的规划口径；日内实际偏差下储能仍可放电至运行下限（1200\,kWh）。

\textbf{（4）日内实际运行与结算。} 日内仍采用与问题二相同的储能实时纠偏机制
（缺额：减充电$\to$加放电$\to$5 倍紧急购电；富余：减放电$\to$加充电$\to$弃光），
并按式~\eqref{eq:q3-cost} 对调整后的实际购电量结算。

\subsection{模型求解与分析}

\subsubsection{参数寻优}
\label{sec:q3-param}
调整机制使预测偏差的代价由 5 倍紧急购电降为 0.5 倍调整费用，
余量参数重新寻优后负载余量由 1.05 降至 $\gamma=1.02$，光伏余量保持 $\delta=0.98$。
当日残差更新权重的寻优结果如表~\ref{tab:q3-lambda} 所示：
不用当日信息（$\lambda=0$）时总费用 1376.3 万元，最优点 $(0.1,0.3)$ 处 1368.8 万元，
说明当日已实现偏差确有信息增量；$\lambda$ 偏大（0.3/0.4 以上）时调整模型追逐残差噪声、
支付多余调整费用，总费用回升；最优点附近曲面平坦，结果对 $\lambda$ 取值不敏感。

\begin{table}[!htbp]
\caption{当日残差更新权重的寻优结果（全年总费用，单位：万元）}\label{tab:q3-lambda}
\centering
\small
\begin{tabular}{lccccccc}
\toprule
$(\lambda_{pv},\lambda_L)$ & $(0,0)$ & $(0.1,0.2)$ & $(0.1,0.3)$ & $(0.15,0.2)$ & $(0.3,0.4)$ & $(0.45,0.6)$ & $(0.6,0.8)$ \\
\midrule
全年总费用 & 1376.3 & 1369.5 & \textbf{1368.8} & 1370.6 & 1377.2 & 1393.0 & 1414.4 \\
\bottomrule
\end{tabular}
\end{table}

日末最低备用约束的取值经变体实验确定：完全自由末端虽比强制当日闭环略省
（1382.6 万元对 1384.3 万元），但储能每天傍晚放空至下限附近、几乎不留备用，紧急购电天数明显增加；
取 $S^{end}_{\min}=4000$\,kWh（约可用容量的 42\%）后费用与稳健性同时改善，
且在 2000$\sim$7000\,kWh 范围内总费用变化平缓（详见第九章）。

\subsubsection{策略变体对比}
为定量回答“是否需要引入其他时刻的预报”，设计一组策略变体从“预报质量”与“调整机制”
两个维度逐步改进，全年滚动仿真结果如表~\ref{tab:q3-variant} 与图~\ref{fig:q3-scheme} 所示。

\begin{table}[!htbp]
\caption{不同购电策略的全年总费用对比（报告期 2025 年 2 月 1 日--12 月 31 日，共 334 天）}\label{tab:q3-variant}
\centering
\small
\setlength{\tabcolsep}{4pt}
\begin{tabular}{lllll}
\toprule
策略 & 说明 & 总费用/万元 & 费用构成/万元 & 相对问题二 \\
\midrule
V0 & 问题二策略（w1 光伏预测、无调整） & 1454.7 & 实付 1367.6 $+$ 紧急 87.0 & --- \\
V1 & 原始 0:00 预报、无调整 & 1809.5 & 实付 1335.7 $+$ 紧急 473.8 & $+$24.4\% \\
V2 & 原始预报 $+$ 双向调整 & 1586.0 & 实付 1458.3 $+$ 调整 110.1 $+$ 紧急 17.7 & $+$9.0\% \\
V3 & 预报标定、无调整 & 1410.7 & 实付 1315.0 $+$ 紧急 95.6 & $-$3.0\% \\
V4D & 标定 $+$ 仅下调调整 & 1407.0 & 实付 1301.4 $+$ 违约 17.5 $+$ 紧急 88.1 & $-$3.3\% \\
V4 & 标定 $+$ 双向调整（强制当日闭环） & 1384.3 & 实付 1312.8 $+$ 调整 33.7 $+$ 紧急 37.8 & $-$4.8\% \\
V4F & 标定 $+$ 双向调整 $+$ 自由末端 & 1382.6 & 实付 1303.6 $+$ 调整 36.3 $+$ 紧急 42.6 & $-$5.0\% \\
\textbf{V4FR} & \textbf{标定 $+$ 双向调整 $+$ 末端备用 $+$ $\lambda$ 寻优} & \textbf{1368.8} & \textbf{实付 1306.2 $+$ 调整 32.1 $+$ 紧急 30.4} & \textbf{$-$5.9\%} \\
--- & 完美预测（理论下界） & 1224.5 & --- & $-$15.8\% \\
\bottomrule
\end{tabular}
\end{table}

\begin{figure}[!htbp]
\centering
\includegraphics[width=0.8\textwidth]{Q3_scheme.pdf}
\caption{各策略变体的全年总费用对比（虚线为完美预测下的理论下界 1224.5 万元）}
\label{fig:q3-scheme}
\end{figure}

由表~\ref{tab:q3-variant} 可知：直接用原始预报替代 w1 预测使总费用由 1454.7 万元升至 1809.5 万元
（$+$24.4\%），说明\textbf{原始预报的偏差远大于 w1 预测}；
但仅依靠调整机制（V2）即可挽回其中 63\% 的损失，
说明\textbf{调整机制本身具有可观的经济价值}——它把 5 倍电价的紧急购电转化为 1.5 倍溢价、
把过配的全部损失削减为 50\% 违约费用。引入预报标定（V3）后总费用降至 1410.7 万元，
说明\textbf{标定是发挥预报价值的前提}；再引入双向调整、末端备用与残差权重寻优（V4FR），
总费用最终降至 1368.8 万元，较问题二策略降低 5.9\%。
对比 V4D 与 V4 可知，仅允许下调会因无法应对光伏低估而多支付约 50 万元紧急购电费
（88.1 万元对 37.8 万元），说明\textbf{双向调整是必要的}。

\subsubsection{求解结果与分析}
最终策略（V4FR）的全年运行指标如表~\ref{tab:q3-annual} 所示。
报告期内实际购电总费用为 \textbf{1368.75 万元}，较问题二策略（1454.65 万元）降低 85.90 万元（$-$5.9\%）；
紧急购电费用 30.40 万元，较问题二的 87.02 万元下降 65.1\%，
紧急购电量由 14.51 万\,kWh 降至 5.59 万\,kWh。
这表明日内滚动调整有效吸收了预测偏差，把原本需 5 倍电价兜底的缺口转化为代价较低的调整购电。

\begin{table}[!htbp]
\caption{问题三最终策略的年度运行指标（报告期 2025 年 2 月 1 日--12 月 31 日）}\label{tab:q3-annual}
\centering
\small
\setlength{\tabcolsep}{5pt}
\begin{tabular}{lrr}
\toprule
指标 & 数值 & 占比 \\
\midrule
计划购电费用（0:00 计划口径） & 12863249.43 元 & --- \\
按调整购电量实付费用 & 13062110.51 元 & 100.0\% \\
调整违约与溢价费用 & 321415.34 元 & 2.5\% \\
紧急购电费用（5 倍电价） & 303964.84 元 & 2.3\% \\
\midrule
全年购电总费用 & \textbf{13687490.69 元（1368.75 万元）} & --- \\
\bottomrule
\end{tabular}
\end{table}

四个指定日期的计划与调整结果如表~\ref{tab:q3-days} 与表~\ref{tab:q3-six} 所示。
调整主要集中在 12:00--12:10 与 18:00--18:10 两个时段：
3 月 20 日 12:00 由计划 629.79\,kWh 上调至 839.81\,kWh；
6 月 21 日 18:00 由 445.00\,kWh 微调至 479.30\,kWh；
9 月 23 日 12:00 由 529.25\,kWh 上调至 852.32\,kWh；
12 月 21 日 12:00 由 996.21\,kWh 上调至 1256.55\,kWh，
16:00 由 830.71\,kWh 上调至 842.74\,kWh，18:00 由 731.57\,kWh 微调至 736.34\,kWh。
四个日期 0:00 与 24:00 的储电量均处于 4115$\sim$5206\,kWh 的健康备用区间，
说明日末备用约束有效避免了储能被放空。

\begin{table}[!htbp]
\caption{四个指定日期的计划购电量、调整后购电量、购电费用与紧急购电情况}\label{tab:q3-days}
\centering
\small
\setlength{\tabcolsep}{5pt}
\begin{tabular}{lcccccc}
\toprule
日期 & 计划全天/kWh & 调整全天/kWh & 计划购电费/元 & 调整后结算/元 & 紧急购电 \\
\midrule
2025-03-20 & 70931.49 & 71641.61 & 43191.79 & 43970.52 & 无 \\
2025-06-21 & 35382.09 & 35651.35 & 20188.83 & 21095.57 & 无 \\
2025-09-23 & 67776.12 & 69583.76 & 42138.53 & 44294.40 & 20:30--20:40，6.38\,kWh \\
2025-12-21 & 96674.20 & 98469.77 & 61965.50 & 63563.48 & 8:00--8:10，29.55\,kWh \\
\bottomrule
\end{tabular}
\end{table}

\begin{table}[!htbp]
\caption{四个指定日期在题目指定时间段内的计划购电量 $\rightarrow$ 调整后购电量（单位：kWh）}\label{tab:q3-six}
\centering
\small
\setlength{\tabcolsep}{4pt}
\begin{tabular}{llcccccc}
\toprule
日期 & 口径 & 10:00--10:10 & 12:00--12:10 & 14:00--14:10 & 16:00--16:10 & 18:00--18:10 & 20:00--20:10 \\
\midrule
2025-03-20 & 计划 & 0.0000 & 629.7876 & 0.0000 & 490.5334 & 737.7354 & 0.0000 \\
           & 调整 & 0.0000 & \textbf{839.8117} & 0.0000 & 490.5334 & 737.7354 & 0.0000 \\
2025-06-21 & 计划 & 0.0000 & 0.0000 & 0.0000 & 123.9504 & 444.9959 & 0.0000 \\
           & 调整 & 0.0000 & 0.0000 & 0.0000 & 123.9504 & \textbf{479.2983} & 0.0000 \\
2025-09-23 & 计划 & 0.0000 & 529.2523 & 0.0000 & 422.6452 & 821.9636 & 0.0000 \\
           & 调整 & 0.0000 & \textbf{852.3236} & 0.0000 & 422.6452 & 821.9636 & 0.0000 \\
2025-12-21 & 计划 & 0.0000 & 996.2066 & 0.0000 & 830.7113 & 731.5661 & 0.0000 \\
           & 调整 & 0.0000 & \textbf{1256.5490} & 0.0000 & \textbf{842.7382} & \textbf{736.3372} & 0.0000 \\
\bottomrule
\end{tabular}
\end{table}

\subsubsection{是否需要引入其他时刻的预报}
综合上述诊断与对比，回答是：\textbf{需要，但必须以正确的方式引入}。

\textbf{（1）后续预报本身的信息增量有限。} 同一目标时刻的预报精度并不随发布时间改善，
其修订量与当日已实现偏差几乎不相关，不宜用最新预报直接替换原计划——
这正是 V1 策略费用比问题二策略高 24.4\% 的原因。

\textbf{（2）后续预报的价值主要体现在“触发调整”与“参与标定”。} 即使原始预报质量不佳，
以其为输入的双向调整机制（V2）仍能挽回 63\% 的预报损失；
经 MOS 标定并与 w1 融合后（V3、V4FR），后续预报才真正转化为经济收益。

\textbf{（3）当日已实现信息是主要的日内信息增量。} 以已实现偏差按适度权重修正后续估计
（$\lambda_{pv}=0.1$、$\lambda_L=0.3$）可使总费用再降 7.5 万元；
权重过大则追逐噪声、增大调整摩擦，残差更新应当保守使用。

\textbf{（4）调整机制的收益来自费率结构的改善。} 0.5 倍的双向调整费率使模型在
“少购违约（$0.5c_t$）”“多购溢价（$1.5c_t$）”与“紧急购电（$5c_t$）”三个通道之间自动比价，
把问题二中被动承担的高成本缺口转化为主动的低成本计划调整。

综上，\textbf{应当引入其他时刻的预报，但应作为经标定与融合后的输入，
并配合当日已实现残差的适度修正与日末备用约束统一纳入滚动调整模型}，
而不是直接以最新预报取代原计划。
